\documentclass[../../../include/open-logic-section]{subfiles}

\begin{document}

\olfileid{sth}{replacement}{absinf}
\olsection{替换与“绝对无穷”}

现在我们将 \limofsize{} 暂且搁置，去探索为替换公理提供辩护的另一类（内在）尝试，这些尝试确实认真对待了集合是分阶段形成的观念。

当我们最初概述迭代过程时，提出了一些解释各个阶段发生何事的原则，即 \stageshier、\stagesord 和 \stagesacc。后来，我们又补充了一些说明阶段数量的原则：\stagessucc{} 表明集合形成的过程永无止境，而 \stagesinf{} 表明该过程会经历一个无限阶段。

我们有理由认为，后两个原则源自某个更宽泛的原则，例如：

\begin{enumerate}
	\item[] \stagesinex. 存在绝对无穷多的阶段；层级的高度达到了其可能达到的极限。
\end{enumerate}

显然，这是一条非形式的原则。但即便它并非由集合的累积迭代观念直接\emph{推出}，它也肯定与之\emph{相协调}。至少，与 \limofsize{} 不同，它保留了集合是逐步形成的想法。

现在，我们希望能借助 \stagesinex{} 为替换公理提供辩护。那么，让我们来看看这是如何做到的。

在 \olref[ordinals][idea]{sec} 中，我们看到很容易构造一个（实质上）理应与 $\omega+\omega$ 同构的良序。换言之，我们很容易想象一个逐步迭代的进程，其序型（实质上）是 $\omega+\omega$。既然如此，如果我们已经接受了 \stagesinex，那么我们当然应该接受层级中至少存在第 $\omega+\omega$ 个阶段，即 $V_{\omega+\omega}$，因为层级显然\emph{可以}延伸得那么远。

这一想法可以推广如下：对于任何良序，构建迭代层级的进程都应至少推进到与该良序同样远。而我们只需将 \olref[ordinals][ordtype]{thmOrdinalRepresentation} 视为一条\emph{公理}，就能保证这一点。它告诉我们，任何良序都同构于某个 冯·诺伊曼序数。由于每个 冯·诺伊曼 序数都等于它自身的秩，\olref[ordinals][ordtype]{thmOrdinalRepresentation} 便告诉我们：每当我们在集合论中能描述一个良序时，集合构建的迭代过程就必定超越该良序。

这显然像是 \stagesinex{} 的一个推论。不幸的是，如果我们的目标是从这个想法中导出替换公理，那么我们就会面临一个简单的技术障碍：替换公理严格强于 \olref[ordinals][ordtype]{thmOrdinalRepresentation}。（这一观察见于 \citet[\S13.2]{Potter2004}；我们将在 \olref[sth][replacement][finiteaxiomatizability]{sec} 中给出证明。）

由此可见，如果我们打算以一种能得出替换公理的方式来理解 \stagesinex{}，那么它就不能\emph{仅仅}断言层级超越了任何良序，而必须做出更强的断言。为此，\cite{Shoenfield:AST} 对该想法提出了一个极其自然的强化，即：层级不与任何集合\emph{共尾}。\footnote{哥德尔 似乎提出过类似的想法；参见 \citet[p.~223]{Potter2004}。关于 哥德尔 和 \citeauthor{Shoenfield:AST} 的讨论，参见 \citet[90--5]{Incurvati2020}。} 更具体地说：如果 $\tau$ 是一个将集合映射到层级阶段的函数，那么任何集合 $A$ 在 $\tau$ 下的像都不能穷尽整个层级。换言之（以模式形式表述）：

\begin{enumerate}
	\item[] \stagescofin. 如果 $A$ 是一个集合，且对每个 $x \in A$，$\tau(x)$ 都是一个阶段，那么存在一个晚于所有 $x \in A$ 对应的 $\tau(x)$ 的阶段。
\end{enumerate}

显然，$\ZF$ 能够证明 \stagescofin{} 经过适当形式化后的版本。反过来，我们可以非形式地论证，\stagescofin{} 为替换公理提供了辩护。\footnote{要用 \stagescofin{} 的某个形式化版本来证明替换公理会更困难，因为单靠 $\Z$ 本身不足以定义阶段，所以不清楚该如何形式化 \stagescofin。不过，一种选择是在 $\LT$ 的某个扩充中展开，正如 \olref[sth][replacement][extrinsic]{sec} 节所述。} 假设 $(\forall x \in A)\lexists![y][\phi(x,y)]$。那么对于每个 $x \in A$，令 $\sigma(x)$ 为使得 $\phi(x,y)$ 成立的那个 $y$，并令 $\tau(x)$ 为 $\sigma(x)$ 首次形成的阶段。根据 \stagescofin，存在一个阶段 $V$，使得 $(\forall x \in A)\tau(x)\in V$。现在，由于每个 $\tau(x) \in V$ 且 $\sigma(x) \subseteq \tau(x)$，由分离公理可得 $\Setabs{y \in V}{(\exists x \in A)\sigma(x) = y} = \Setabs{y}{(\exists x \in A)\phi(x,y)}$。

\begin{prob}
	在 $\ZF$ 内形式化 \stagescofin{}。
\end{prob}

所以，\stagescofin{} 为替换公理提供了辩护。并且，至少有理由认为 \stagesinex{} 也为 \stagescofin{} 提供了辩护。假设 \stagescofin{} 不成立，那么该层级相对于某个集合~$A$ 是共尾的，即我们有一个映射 $\tau$，使得对任意阶段~$S$，都存在某个 $x \in A$ 使得 $S \in \tau(x)$。在这种情况下，我们确实有办法把握该层级所谓的“绝对无限性”：它被 $\tau$ 作用于 $A$ 的\emph{值域}所穷尽。这就削弱了“该层级是绝对无限的”这一想法。换质位之：\stagesinex{} 语义地后承 \stagescofin，而后者又为替换公理提供了辩护。

这确实是一次真正有希望的尝试，旨在为替换公理提供内在辩护。但它最终是否成立，只能留待你来判断。

\end{document}